\documentclass[english, parskip=full]{scrartcl}
\usepackage[utf8]{inputenc}  % Kodierung der Datei
\usepackage[english]{babel}  % Sprache auf Deutsch 
\usepackage{color}
\usepackage{graphicx, gensymb}
\usepackage{amsfonts, cancel, amssymb}
\usepackage{amsmath, amsthm, array, esint}
\usepackage{booktabs, multirow}  % schönere Tabellen
\usepackage{caption}  % Anpassung der Captions
\usepackage{scrlayer-scrpage}    % Manipulation des Seitenstils
%\pagestyle{scrheadings}  % Stil für die Seite setzen
%\automark{section}  % Abschnittsnamen als Seitenbeschriftung 
%\setheadsepline{.5pt}    % Kopzeile durch Linie abtrennen
\usepackage[hidelinks]{hyperref}  % Links und weitere PDF-Features
\usepackage{typearea, tikz, pgffor, verbatim}
\usetikzlibrary{calc, quotes}
\usepackage[cache=false, outputdir=build]{minted}
\usepackage{placeins} % provides \FloatBarrier

\definecolor{bg}{rgb}{0.9,0.9,0.9}
\newcommand\numberthis{\addtocounter{equation}{1}\tag{\theequation}}
\newcommand{\bra}{}
\newcommand{\kom}{}
\newcommand{\brak}{}
\newcommand{\braket}{}
\newcommand{\intx}{}
\newcommand{\intr}{}
\newcommand{\kla}{}
\newcommand{\klam}{}
\newcommand{\klammer}{}
\newcommand{\oper}{}
\newcommand{\tecxt}{}
\newcommand{\Bra}{}
\newcommand{\Ket}{}
\renewcommand{\i}{\text{i}}
\newcommand{\e}{\text{e}}
\newcommand*{\Fourier}{\textsc{Fourier}\ }
\newcommand*{\F}{\mathcal{F}}
\newcommand*{\sinc}{\text{sinc\,}}
\newcommand{\conv}{\mathop{\scalebox{3.5}{\raisebox{-0.38ex}{\makebox[0.5\width][c]{$\ast$}}}}\limits}

\newcommand*{\titel}{03 Quadratur}
\newcommand*{\autor}{Joachim Schwardt und Julian Fleck}

\hypersetup{pdfauthor={\autor}, pdftitle={\titel}}

%\titlehead{titlehead}
\subject{Numerik I - Aufgaben}
\title{\titel}
%\author{\autor}
\author{\begin{tabular}{l|l}
Joachim Schwardt & 4768711 \\ 
Julian Fleck & 4759587\\ 
\end{tabular}}
%\author{\begin{tabular}{|l|l|}
%\hline
%Joachim Schwardt & 4768711 \\ \hline
%Julian Fleck & 4759587\\ \hline
%\end{tabular}}
\date{\today}

\begin{document}

%\begin{titlepage}
\maketitle  % Titel setzen
%\tableofcontents  % Inhaltsverzeichnis setzen
%\end{titlepage}

\section*{Aufgabe 1}
Gegeben ist die Quadraturformel
\begin{align}
Q_{[a,b]}[f]&:= \sum_{j=1}^{n}\frac{h_j}{2}\klam\left[ f(x_j) + f(x_{j-1}) \right]+\frac{h_j^2}{12}\klam\left[ f'(x_{j-1}) - f'(x_{j}) \right].
\end{align}
\subsection*{a)}
Das exakte Integral eines kubischen Spline $s_3$ mit den Stützstellen $a=x_0<\dots<x_n=b$ ist gegeben durch
\begin{align}
J[s_3] &:= \intx\int_{a}^{b}s_3(x)\, \mathrm{d}x = \\
&= \sum_{j=0}^{n-1} \intx\int_{x_j}^{x_{j+1}} a_j(x-x_j)^3 + b_j(x-x_j)^2 + c_j(x-x_j) + d_j \, \mathrm{d}x = \\
&= \sum_{j=0}^{n-1}\frac{a_j}{4}(x_{j+1}-x_j)^4 + \frac{b_j}{3}(x_{j+1}-x_j)^3 + \frac{c_j}{2}(x_{j+1}-x_j)^2 + d_j(x_{j+1}-x_j) = \\
&= \sum_{j=1}^{n}\frac{a_{j-1}}{4}h_j^4 + \frac{b_{j-1}}{3}h_j^3 + \frac{c_{j-1}}{2}h_j^2 + d_{j-1}h_j .
\end{align}
Die Quadraturformel ergibt dabei
\begin{align*}
Q_{[a,b]}[s_3] &=\sum_{j=1}^n \frac{h_j}{2}\klam\left[ d_{j-1} + a_{j-1}h_j^3 + b_{j-1}h_j^2 + c_{j-1}h_j + d_{j-1} \right] +\\
&\hspace{1cm} + \frac{h_j^2}{12}\klam\left[ c_{j-1} - 3a_{j-1}h_j^2 - 2b_{j-1}h_j - c_{j-1} \right] = \numberthis\\
&=\sum_{j=1}^n a_{j-1}h_j^4\kla\left( \frac{1}{2} - \frac{3}{12} \right)+b_{j-1}h_j^3\kla\left( \frac{1}{2} - \frac{2}{12} \right) + c_{j-1}h_j^2\kla\left( \frac{1}{2} \right)+d_{j-1}h_j= \numberthis\\
&=\sum_{j=1}^n\frac{a_{j-1}}{4}h_j^4 + \frac{b_{j-1}}{3}h_j^3 + \frac{c_{j-1}}{2}h_j^2 + d_{j-1}h_j\equiv J[s_3]. \numberthis
\end{align*}
Kubische Splines können mit der Quadraturformel also exakt integriert werden.
\subsection*{b)}
Für $h_j\equiv\tecxt\text{const.}=:h$ kann die Quadraturformel vereinfacht werden, da
\begin{align}
Q[f] &= \frac{h}{2}\sum_{j=1}^n f(x_j) + \frac{h}{2}\sum_{j=1}^n f(x_{j-1}) + \frac{h^2}{12}\sum_{j=1}^nf'(x_{j-1}) - \frac{h^2}{12}\sum_{j=1}^n f'(x_{j})=\\
&= \frac{h}{2}\sum_{j=1}^n f(x_j) + \frac{h}{2}\sum_{j=0}^{n-1} f(x_{j}) + \frac{h^2}{12}\sum_{j=0}^{n-1} f'(x_{j}) - \frac{h^2}{12}\sum_{j=1}^nf'(x_j)=\\
&= \frac{h}{2}\kla\left( f(x_0) + f(x_n) \right)+h\sum_{j=1}^{n-1} f(x_j) + \frac{h^2}{12}\kla\left( f'(x_0) - f'(x_n) \right).
\end{align}

\newpage 
\section*{Aufgabe 2}
Gegeben ist die Quadraturformel
\begin{align}
Q_n[f] &:= \sum_{j=1}^n a_jf(x_j,y_j),
\end{align}
um Integrale von Funktionen auf dem Einheitsdreieck $T:=\klammer\left\{ (0,0),(0,1),(1,0) \right\}$ zu approximieren.
\subsection*{a)}
Sei $n=1$ und die (einzige) Stützstelle gegeben durch den geometrischen Schwerpunkt des Dreiecks, d.h.
\begin{align}
x_1&=y_1 = \frac{1}{3}.
\end{align}
Die Quadraturformel lautet also explizit
\begin{align}
Q_1[f]&=a_1f\kla\left( \frac{1}{3},\frac{1}{3} \right).
\end{align}
Sei $p_1(x,y):=ax+by+c$ ein beliebiges Polynom vom Grad eins.
Dann ist
\begin{align}
J[p_1]&:=\intx\int_{T}^{}p_1(x,y)\, \mathrm{d}(x,y)=\intx\int_{0}^{1}\intx\int_{0}^{1-x}ax+by+c\, \mathrm{d}y\, \mathrm{d}x=\\
&= \intx\int_{0}^{1}ax(1-x)+c(1-x) + \frac{b}{2}(1-x)^2\, \mathrm{d}x =\\
&= \kla\left( \frac{x^3}{3}\klam\left[ -a+\frac{b}{2} \right]+\frac{x^2}{2}\klam\left[ a-b-c \right]+x\klam\left[ c + \frac{b}{2} \right] \right)\bigg\vert_0^1 =\\
&=\frac{a}{6}+\frac{b}{6}+\frac{c}{2}.
\end{align}
Andererseits ist
\begin{align}
Q_1[p_1] &= a_1\kla\left( \frac{a}{3} + \frac{b}{3} + c \right),
\end{align}
für $a_1:=\frac{1}{2}$ folgt daher $J[p_1]\equiv Q_1[p_1]$.
Der minimale Exaktheitsgrad ist damit also gleich eins.
Jetzt sind allerdings alle Parameter festgelegt, der maximale Exaktheitsgrad ist somit ebenfalls eins.
Betrachte dazu einfach das Polynom $p_2(x,y)=ax^2$.
Dann ist
\begin{align}
J[p_2] &= \intx\int_{0}^{1}\intx\int_{0}^{1-x}p_2(x,y)\, \mathrm{d}y\, \mathrm{d}x = \intx\int_{0}^{1}ax^2(1-x)\, \mathrm{d}x = a\kla\left( \frac{1}{3}-\frac{1}{4} \right)=\frac{a}{12}.
\end{align}
Allerdings liefert die Quadraturformel hier
\begin{align}
Q_1[p_2] &= \frac{1}{2}\cdot a\kla\left( \frac{1}{3} \right)^2 = \frac{a}{18}\neq J[p_2].
\end{align}
\subsection*{b)}
Sei nun $n=3$ und die (drei) Stützstellen gegeben durch die Mittelpunkte der Kanten des Dreiecks $T$, also
\begin{align}
(x_1,y_1)&=\kla\left( 0,\frac{1}{2} \right),&&&\kla\left( x_2,y_2 \right)&=\kla\left( \frac{1}{2},0 \right),&&&\kla\left( x_2,y_2 \right)&= \kla\left( \frac{1}{2},\frac{1}{2} \right).
\end{align}
Die Quadraturformel lautet also explizit
\begin{align}
Q_2[f]&= a_1f\kla\left( 0,\frac{1}{2} \right) + a_2f\kla\left( \frac{1}{2}, 0 \right) + a_3f\kla\left( \frac{1}{2},\frac{1}{2} \right).
\end{align}
Sei $p_2(x,y):=ax^2+bxy+cy^2+dx+ey+f$ ein beliebiges Polynom vom Grad zwei.
Dann ist
\begin{align*}
J[p_2]&:=\intx\int_{T}^{}p_2(x,y)\, \mathrm{d}(x,y)=\intx\int_{0}^{1}\intx\int_{0}^{1-x}ax^2+bxy+cy^2+dx+ey+f\, \mathrm{d}y\, \mathrm{d}x= \numberthis\\
&= \intx\int_{0}^{1}ax^2(1-x)+\frac{b}{2}x(1-x)^2 + \frac{c}{3}(1-x)^3 +\\
&\hspace{1cm} + dx(1-x) + \frac{e}{2}(1-x)^2 + f(1-x)\, \mathrm{d}x = \numberthis\\
&= \kla\bigg( \frac{x^4}{4}\klam\left[ -a + \frac{b}{2} - \frac{c}{3} \right] + \frac{x^3}{3}\klam\left[ a-b+c - d+\frac{e}{2} \right]+ \\
&\hspace{0.5cm} + \frac{x^2}{2}\klam\left[ \frac{b}{2} -c +d -e -f \right]•  +x\klam\left[ \frac{c}{3} + \frac{e}{2} + f \right] \bigg)\bigg\vert_0^1 =\numberthis\\
&= \frac{a}{12} + \frac{b}{24} +\frac{c}{12} + \frac{d}{6} + \frac{e}{6} +\frac{f}{2} \numberthis .
\end{align*}
Andererseits ist
\begin{align*}
Q_2[p_2] &= a_1\kla\left( \frac{c}{4}+\frac{e}{2}+f \right) + a_2\kla\left( \frac{a}{4} + \frac{d}{2} + f \right) + a_3\kla\left( \frac{a+b+c}{4}+\frac{d+e}{2}+f \right) = \numberthis \\
&=a\kla\left( \frac{a_2+a_3}{4} \right)+b\kla\left( \frac{a_3}{4} \right)+c\kla\left( \frac{a_1+a_3}{4} \right) +\\
&\hspace{0.5cm} + d\kla\left( \frac{a_2+a_3}{2} \right)+ e\kla\left( \frac{a_1 + a_3}{2} \right)+f\kla\left( a_1+a_2+a_3 \right) \numberthis  .
\end{align*}
Wir erhalten somit die Bedingungen $a_3=\frac{1}{6}$, $a_2+a_3 = \frac{1}{3}$, $a_1+a_3=\frac{1}{3}$ und $a_1+a_2+a_3 = \frac{1}{2}$.
Man aus der zweiten und dritten Bedingung direkt $a_1=a_2=\frac{1}{3}-a_3=\frac{1}{6}$.
Die letzte Bedingung ist tatsächlich auch erfüllt, da $a_1+a_2+a_3=\frac{1+1+1}{6}=\frac{1}{2}$.
Der minimale Exaktheitsgrad ist also hier gleich zwei.
Genau wie zuvor haben wir allerdings bereits alle Parameter festgelegt, weshalb leicht zu zeigen ist, dass der Exaktheitsgrad auch maximal gleich zwei ist.
Betrachte dazu einfach das Polynom $p_3(x,y)=ax^3$.
Dann ist
\begin{align}
J[p_3] &= \intx\int_{0}^{1}\intx\int_{0}^{1-x}p_3(x,y)\, \mathrm{d}y\, \mathrm{d}x = \intx\int_{0}^{1}ax^3(1-x)\, \mathrm{d}x = a\kla\left( \frac{1}{4}-\frac{1}{5} \right)=\frac{a}{20}.
\end{align}
Allerdings liefert die Quadraturformel hier
\begin{align}
Q_2[p_3] &= \frac{1}{6}\cdot a\kla\left( \frac{1}{2} \right)^3 + \frac{1}{6}\cdot 0 + \frac{1}{6}\cdot a\kla\left( \frac{1}{2} \right)^3 = \frac{a}{24}\neq J[p_3].
\end{align}


\newpage
\section*{Aufgabe 3}
Gegeben ist die Quadraturformel
\begin{align}
Q_{n}[f]&:= (b-a)\sum_{j=0}^{n}a_jf(x_j).
\end{align}
Sei $q\in\Pi_{n+1}$ ein Polynom vom Grad $n+1$ und gegeben durch 
\begin{align}
q(x) &:=\prod_{j=0}^n \kla\left( x - x_j \right).
\end{align}
Dann folgt offensichtlich
\begin{align}
Q_n[q] &= (b-a)\sum_{j=0}^n a_jq(x_j)  = 0,
\end{align}
da per Konstruktion die Stützstellen $x_j$ per Konstruktion Nullstellen von $q$ sind.
Betrachte nun ein Polynom vom Grad $2n+2$, welches über $p(x):=q(x)^2$ definiert sei.
Dann sind die Stützstellen natürlich auch Nullstellen von $p$, und somit ist $Q_n(p)=0$.
Allerdings gilt per Konstruktion $p(x)\ge 0$, weshalb
\begin{align}
I[p] &:= \intx\int_{a}^{b}p(x)\, \mathrm{d}x\ge 0
\end{align}
folgt.
%Insbesondere kann das Integral über eine nicht-negative Funktion (wie $p(x)$) nur genau dann gleich null sein, wenn die Funktion (fast) überall verschwindet (fast im Sinne von Maßtheorie).
%Das Integral ist also nicht gleich null, womit die Quadraturformel für $p$ nicht exakt ist.
%\\
%Um Maßtheorie zu vermeiden kann aber auch elementar abgeschätzt werden, dass das Integral streng größer als null ist.
Man kann nun abschätzen, dass das Integral streng größer als null ist.
Betrachte dazu einfach
\begin{align}
I[p]&\ge\intx\int_{a}^{x_1}p(x)\, \mathrm{d}x = \intx\int_{a}^{x_1}\prod_{j=0}^n\kla\left( x-x_j \right)^2\, \mathrm{d}x \ge \\
&\ge \underbrace{\prod_{j=2}^n \kla\left( x_j - x_1 \right)^2}_{=:\, C\ > \ 0}  \intx\int_{a}^{x_1}\kla\left( x - a \right)^2\kla\left( x_1 - x \right)^2\, \mathrm{d}x \kom\overset{\substack{{y\,:=\,x-a\ \tecxt\text{und } t\,:=\,x_1 - a} \\ |}}{=} \\
&= C\intx\int_{0}^{t}y^2\kla\left( t - y \right)^2\, \mathrm{d}y = C\kla\left( \frac{y^5}{5} + \frac{y^3t^2}{3} - \frac{ty^4}{2} \right)\bigg\vert_{\cancel{0}}^t = \\
&= Ct^5\kla\left( \frac{1}{5}+\frac{1}{3}-\frac{1}{2} \right)=Ct^5\kla\left( \frac{1}{5} - \frac{1}{6} \right) > 0.
\end{align}
Das Integral ist also nicht gleich null, womit die Quadraturformel für $p$ nicht exakt ist.
Ersteres ist natürlich ``offensichtlich'', weil das Integral über eine beliebige nicht-negative Funktion nur dann verschwinden kann, wenn die Funktion (fast) überall gleich null ist.
Um aber mögliche Fehler bei einer solchen ``hand-wavy'' Begründung zu vermeiden, besser die kurze Rechnung.



%\begin{thebibliography}{9}
%\bibitem{example} description, \url{https://www.exmapleurl}, day.\,month~2021.
%\end{thebibliography}

\end{document}